\documentclass{newlayout}
%Bitte hier den enstprechenden Ort einsetzen z.B. Braunschweig und die Akademienummer
\Akademie{Braunschweig}{2015}{1}

\usepackage[english, ngerman]{babel}
\usepackage{misc}
\usepackage{multicol}

%\usepackage{amsmath}%wird automatisch durch newlayout.cls geladen
\usepackage{amsfonts}
\usepackage{blindtext}

\usepackage{url}
\def\UrlBreaks{\do\a\do\b\do\c\do\d\do\e\do\f\do\g\do\h\do\i\do\j\do\k\do\l%
\do\m\do\n\do\o\do\p\do\q\do\r\do\s\do\t\do\u\do\v\do\w\do\x\do\y\do\z\do\0%
\do\1\do\2\do\3\do\4\do\5\do\6\do\7\do\8\do\9\do\-\do\_\do\/\do\%}
\urlstyle{same}

% hinzugef�gt, um Fehler 'pdfTeX error (font expansion): auto expansion is only possible with scalable' zu vermeiden
\usepackage{lmodern}
\setkomafont{descriptionlabel}{\normalfont\bfseries}
\addtokomafont{paragraph}{\normalfont}
\usepackage{footnote}
\usepackage[flushmargin,hang,ragged]{footmisc}
\deffootnote{1em}{1em}{%
\textsuperscript{\thefootnotemark\ }
}
%\setlength{\abovedisplayskip}{5pt}
%\setlength{\belowdisplayskip}{5pt}


%%%%%Mathe-Definitionen
\newtheorem{Def}{Definition}
\newtheorem{Sat}{Satz}
\newtheorem{Bew}{Beweis}
\newtheorem{Thm}{Theorem}

\setlength\abovedisplayshortskip{0pt}
\setlength\belowdisplayshortskip{0pt}
\setlength\abovedisplayskip{3pt}
\setlength\belowdisplayskip{3pt}
%%%%Ende Mathe-Definitionen

\begin{document}

 %   \input{titel}
 \setcounter{page}{3}

\setcounter{tocdepth}{1}
 \tableofcontents

   \setcounter{secnumdepth}{1}


\setcounter{page}{7}
\setcounter{chapter}{0}

%Angabe, bis zu welcher Stufe die sections im Text nummeriert werden sollen.
      \settocdepth{2}



\course{1}{Projekte}%%% 
\begin{coursetitle}
  \centerline{Schrifterkennung}
  \bigskip
  \Large \centerline{mit SVMs und Neuronalen Netzen}
  \bigskip
 \label{fig:meinbild}
  \bigskip
\end{coursetitle}
%\begin{dsafigure}
%\begin{center}
%\includegraphics[width=.9\columnwidth]{Titelbild-fehlt.png}
%\caption{meine Bildunterschrift}
%label{fig:meinbild}
%\end{center}
%\end{dsafigure}




\section{Schrifterkennung mit SVMs}

\paragraph{Einführung}
Täglich werden in Deutschland rund 64 Millionen Briefe verschickt. Doch habt ihr euch schonmal gefragt, wie die ganzen Briefe so schnell verarbeitet werden können? Die Vorstellung von mehreren tausend Mitarbeitern, die sich nur mit der Sortierung der Post beschäftigen, mag eine einfache Option darstellen, jedoch wäre sie viel zu mühsam. Eine bessere Variante stellt die Benutzung von Schrifterkennungsalgorithmen dar, deren sich die Post seit längerer Zeit bedient. Dabei scannt ein Computer die auf den Umschlägen geschriebenen Adressen ein, wertet diese aus und sortiert sie entsprechend. In unserem Projekt haben wir uns mit einem solchen Algorithmus beschäftigt und diesen mithilfe von SVMs implementiert. Wie das Ganze funktioniert wollen wir im folgenden erläutern.

\paragraph{Aufbau des Programms}

Da es beim maschinellen Lernen essenziell ist über ein entsprechend großes Datenset zu verfügen, haben wir unseren Algorithmus zunächst einmal mit 60.000 Bildern ausgestattet. 

\begin{verbatim}
	feature, hog_image = feats.hog(sub_pic, orientations = 8, pixels_per_cell = (4, 4), cells_per_block = (1, 1), visualise = true)
\end{verbatim}
Im ersten Schritt generieren wir die Features sämtlicher Bilder. Dabei gibt es mehrere Einstellungsmöglichkeiten. Beispielsweise zur Variation der Pixelzahl in einer Zelle und der Genauigkeit zur Erkennung von Konturen. Dabei gilt: Je kleiner die Zellen sind und je höher die Genauigkeit (Orientations) ist, desto länger braucht der Algorithmus zum berechnen, aber desto präziser ist das Endergebnis. Hierbei müssen wir beachten, dass die Zellengröße nicht zu klein wird (min 2x2 Pixel)!
Anschließend wird die SVM mit den Bildern trainiert und ist somit einsatzbereit.Um nun die Genauigkeit der SVM zu testen, erzeugen wir ein Testset, dessen Features ebenfalls über den bestehenden Algorithmus generiert werden. Getetstet wird die SVM, indem wir das Testset an die SVM übergeben und sie uns aufgrund der vorhandenen Werte eine Schätzung abgibt, was am wahrscheinlichsten auf dem Bild abgebildet ist. Die Ergebnisse werden dann überprüft und die Genauigkeit wird errechnet.
Daraufhin erzeugen wir ein Bild mit 100 Zahlen, die weitgehend zufällig auf dem Bild verteilt sind. Dieses Bild wird in eine Matrix umgewandelt. Ein Algorithmus iterriert dann über das gesamte Bild und die SVM ermittelt welche 100 Zahlen abgebildet sind.




\section{Schrifterkennung mit neuronalen Netzen}

Ein funktionsfähiges neuronales Netz zu implementieren bedeutet einen großen Zeitaufwand und viel Feinoptimierung. Da uns diese Zeit nicht zur Verfügung stand und wir ein größeres Interesse an der Anwendung hatten, nutzten wir Caffe. Bei Caffe handelt es sich um eine Software des Berkeley Vision and Learning Centers. Caffe bietet eine Art Bausatz für neuronale Netze. Man kann verschiedene Filter auswählen und sich so sein eigenes neuronales Netz erstellen. Caffe ist in der Lage Bilder sehr schnell auszuwerten, nach Herstellerangaben benötig es je Bild 4 ms in der Trainingsphase und 1 ms in der Testphase.

%\input{Kapitel-2-Kapitel�berschrift}
%\input{Kapitel-3-Kapitel�berschrift}
%usw.

%\nocite{*} % alle referenzen anzeigen, sogar wenn sie nicht im text zitiert sind
\bibliography{lit}{}
\bibliographystyle{plain}
\end{document}
